\section{Twistor spaces and the quaternionic K\"ahler metric}
\label{tw:section}

To prove Corollary~\ref{thm:quaternionic}, we apply
Theorem~\ref{thm:main} to the twistor space of a positive quaternionic
K\"ahler manifold. The complex classification then determines
the Riemannian metric after its scale has been fixed.

\subsection{The normalized metric and its twistor space}

For $m\geq2$, the holonomy containment in the corollary makes $g$
Einstein \cite{Salamon1982}; see also \cite[Definition~1.1 and the
following discussion]{LeBrun1995}. Thus, writing $s>0$ for its constant
scalar curvature, we have $\operatorname{Ric}_g=(s/4m)g$.
The manifold $M$ is compact and connected by hypothesis.

We fix the scale at this point. Set
\begin{equation}\label{tw:normalization}
\widehat g=\frac{s}{16m(m+2)}\,g,
\qquad
\operatorname{Scal}_{\widehat g}=16m(m+2).
\end{equation}
A constant positive rescaling preserves the Levi--Civita connection
and the parallel quaternionic structure.

The holonomy reduction determines a parallel rank-three bundle
$Q\subset\End(TM)$ whose fibers are the imaginary
quaternions. Its sphere of complex structures is the twistor bundle
\[
\pi:X=\{J\in Q:J^2=-\id\}\longrightarrow M.
\]
The Levi--Civita connection gives a horizontal distribution. Combining
the complex structure $J$ on each horizontal tangent space with the
standard complex structure on the sphere fiber gives the twistor
complex structure. The positive twistor theorem
\cite{Salamon1982}, in the form of
\cite[Theorem~2.1]{LeBrun1995}, says that it is integrable and that
$X$ is a complex contact manifold of dimension $2m+1$. Its contact
distribution $F$ is horizontal, and $X$ carries a K\"ahler--Einstein
metric of scalar curvature $8(m+1)(2m+1)$.
The fibers are holomorphically embedded copies of $\PP^1$ with normal
bundle $\OO(1)^{\oplus2m}$. In the normalized metrics, $\pi$ is a
Riemannian submersion with totally geodesic sphere fibers of curvature
$4$. The antipodal map on the fibers is a free antiholomorphic
involution of $X$.

The total space $X$ is compact and connected, since it is a sphere
bundle over the compact connected manifold $M$.
The positive Ricci form of its K\"ahler--Einstein metric is the
curvature form of the anticanonical bundle. The Kodaira embedding
theorem \cite{Kodaira1954} therefore makes $X$ projective with
$-K_X$ ample. The holomorphic contact sequence is algebraic by
GAGA \cite{Serre1956}. Writing $L=T_X/F$, the contact pairing gives
\begin{equation}\label{tw:determinants}
\det F\simeq L^{\otimes m},
\qquad -K_X\simeq L^{\otimes(m+1)}.
\end{equation}
Thus $L$ is ample as well.

The Fano property also settles the topology of the original manifold.
Lemma~\ref{red:simply-connected} gives $\pi_1(X)=0$. Since the fiber
$S^2$ is connected and simply connected, the homotopy exact sequence of
$X\to M$ gives
\[
\pi_1(M)\simeq\pi_1(X)=0.
\]
Thus the metric recovery below will identify $M$ itself with the
symmetric model.

The cohomological properties used in the algebraic argument follow here
as well. Kodaira vanishing
\cite{Kodaira1953}, applied to $K_X\otimes(-K_X)$, gives
$H^1(X,\OO_X)=H^2(X,\OO_X)=0$. The exponential sequence identifies
$\Pic(X)$ with $H^2(X,\ZZ)$, which is finitely generated and
torsion-free by the universal coefficient theorem and
$H_1(X,\ZZ)=0$.

Every point of $X$ lies on its own twistor fiber. For a parametrization
$u:\PP^1\hookrightarrow X$ of any such fiber, horizontality makes
$u^*F$ complementary to $du(T_{\PP^1})$. The normal bundle assertion
and the contact quotient $\theta:T_X\to L$ consequently give
\begin{equation}\label{tw:fiber-splitting}
\begin{gathered}
u^*F\simeq N_{u(\PP^1)/X}\simeq\OO(1)^{\oplus2m},
\qquad
\theta\,du:T_{\PP^1}\xrightarrow{\sim}u^*L,\\
u^*L\simeq\OO(2),\qquad
u^*T_X=du(T_{\PP^1})\oplus u^*F
\simeq\OO(2)\oplus\OO(1)^{\oplus2m}.
\end{gathered}
\end{equation}
In particular, these are contact-degree-two curves transverse to $F$
through every point.

\subsection{The symmetric model and metric recovery}

\begin{proof}[Proof of Corollary~\ref{thm:quaternionic}]
The preceding construction produces a smooth connected projective
contact Fano manifold $(X,F)$. Theorem~\ref{thm:main} therefore gives
a contact isomorphism from $(X,F)$ to the adjoint contact variety
of a simple complex Lie algebra. It also identifies $L$ with the
restriction of the hyperplane bundle in the adjoint embedding.

Wolf's correspondence \cite{Wolf1965} identifies this adjoint variety
with the complex twistor space of a compact simply connected symmetric
positive quaternionic K\"ahler model $(M_0,g_0)$; the correspondence
for every simple Lie type is described in
\cite[Section~3.9, pp.~105--106]{LandsbergManivel2004}.
Since the twistor space has complex dimension $2m+1$, the model has
real dimension $4m$. Rescale $g_0$ to a metric $\widehat g_0$ of scalar
curvature $16m(m+2)$, as in \eqref{tw:normalization}.

LeBrun's metric uniqueness theorem
\cite[Corollary~3.4]{LeBrun1995}, with the standing normalization
specified after Definition~1.2 of that paper, now applies: the
biholomorphism of the two complex twistor spaces yields an isometry
\[
(M,\widehat g)\simeq(M_0,\widehat g_0).
\]
The corollary asserts the existence of a base isometry from the
complex isomorphism class of the twistor space. It does not require
the initially given biholomorphism to preserve the chosen real
structures or twistor fibrations.

Undoing the two constant rescalings shows that $(M,g)$ itself is
homothetic to the compact symmetric Wolf space $(M_0,g_0)$. Its curvature
is therefore parallel. Indeed, constant rescaling leaves the
Levi--Civita connection unchanged and multiplies the covariant curvature
tensor by the same constant, so it preserves the equation
$\nabla\operatorname{Rm}=0$. Hence
\[
\nabla\operatorname{Rm}_g=0.
\]
\end{proof}
